%% 24_ecg.tex -- Chapter 24: ECG Processing
%% Status: OUTLINE (build stage 5). Replace the outline with the full chapter when it is written.
\chapter{ECG Processing}\label{ch:ecg}
\chapterlead{Detecting heartbeats is the oldest biomedical DSP problem, and the best way to see how ADC bits and amplifier noise reach a clinical result.}

\begin{specnote}[Chapter status: outline]
This chapter is an outline. It will be written in build stage 5 (Part IV, including bioimpedance, with Labs 11--13 and 15). The topics, the design question and the labs below are fixed; the derivations, worked examples and figures are still to come.
\end{specnote}

\begin{objectives}
\begin{itemize}
  \item implement the Pan--Tompkins QRS detector;
  \item compute heart rate and heart-rate variability metrics;
  \item classify beats on the MIT-BIH Arrhythmia Database;
  \item test robustness with the MIT-BIH Noise Stress Test Database;
  \item show how ADC bits and AFE noise move detection sensitivity.
\end{itemize}
\end{objectives}

\begin{insight}[Design question]
How many bits and how much input noise can the ECG chain afford before QRS detection degrades?
\end{insight}

\section{Pan--Tompkins QRS detection}\label{sec:ch24-1}
Band-pass, derivative, squaring, moving-window integration and adaptive thresholds~\cite{pantompkins1985}.

\section{Heart rate and HRV}\label{sec:ch24-2}
RR intervals, time-domain and frequency-domain HRV metrics.

\section{Beat classification}\label{sec:ch24-3}
Features and simple classifiers on the MIT-BIH Arrhythmia Database~\cite{mitbih}.

\section{Noise stress testing}\label{sec:ch24-4}
Calibrated baseline wander, muscle and electrode-motion noise from the Noise Stress Test Database~\cite{nstdb}.

\section{Co-design: bits and noise versus sensitivity}\label{sec:ch24-5}
Sweeping the front-end parameters of a behavioural chain against detection sensitivity and positive predictivity.

\begin{tryit}[Labs for this chapter]
Lab~11 (\cref{sec:lab11}) passes MIT-BIH and Noise Stress Test records through a behavioural AFE and ADC and plots QRS sensitivity and PPV versus ADC bits and input noise.
\end{tryit}

\paragraph{Read with this chapter} Pan and Tompkins~\cite{pantompkins1985}, Tompkins~\cite{tompkins}, Moody and Mark~\cite{mitbih}.
